%!TEX Root = ../main.tex
\chapter{Background}
\label{chap:background}
This appendix provides some background material.
\Cref{sec:category} and \Cref{sec:universal} 
introduces basic category theory
(readers interested in learning further are encouraged
to consulte \cite{category}).
\Cref{sec:topological_space} gives a quick review 
of topological spaces and topological notions.

\section{Category and Functors}
\label{sec:category}
Any mathematical structure you will care about
forms a category.
\begin{definition}
    A \vocab{category} $\mathcal C$ consists of:
    \begin{itemize}
    \item a set of \vocab{objects}
    $\operatorname{obj} \mathcal C$;
    \notation{$\operatorname{obj}\calC$}{set of objects}
    \item  for any
    $A,B\in\operatorname{obj} \mathcal C$,
    the set of \vocab{morphisms} $\operatorname{Mor}(A,B)$; and
    \notation{$\operatorname{Mor}(A,B)$}{set of morphisms}
    \item for any $A,B,C\in\operatorname{obj} \mathcal C$,
    the composition operation $\circ$
    which takes morphisms, one in $\operatorname{Mor}(A,B)$
    and one in $\operatorname{Mor}(B,C)$
    to get a morphism $\operatorname{Mor}(A,C)$
    \end{itemize}
    satisfying the following conditions 
    (which are mild)
    \begin{itemize}
    \item \textbf{associativity.} 
    $f\circ(g\circ h) = (f\circ g)\circ h$
    for any compatible morphisms $f,g,h$.
    \item \textbf{identity.}
    for any $A\in \operatorname{obj}(\mathcal C)$,
    there exists identity morphism $1_A \in 
    \operatorname{Mor}(A,A)$ such that 
    $1_A\circ f = f\circ 1_B = f$
    for all $f\in \operatorname{Mor}(A,B)$.
    \end{itemize}
\end{definition}
\begin{example}
\notation{$\Rings$}{category of rings}
Commutative rings form a category
$\Rings$ with objects being rings 
and morphisms being ring homomorphisms.

Similarly, the category $\Mod R$ of $R$-modules
has objects being $R$-modules and morphisms being 
$R$-module homomorphisms.
A special case of this is vector spaces (when $R=k$ is a field)
and abelian groups (when $R=\ZZ$).
\end{example}
We can also study maps between two categories,
called the \vocab{functor}.
\begin{definition}
    A \vocab{covariant functor} from $\mathcal C\to\mathcal D$
    consists of two maps:
    $F^{\text{obj}} : \operatorname{obj} \mathcal C
    \to\operatorname{obj} \mathcal D$
    and $F^{\text{Mor}} : 
    \operatorname{Mor}(A,B) \to \operatorname{Mor}(F(A),F(B))$
    satisfying the following conditions:
    \begin{itemize}
    \item \textbf{preserves composition.}
    $F(f\circ g) = F(f)\circ F(g)$.
    \item \textbf{perserves identity.}
    $F(1_A) = 1_{F(A)}$ for all $A\in\operatorname{obj}\calC$
    \end{itemize}

    (In contrast, a \vocab{contravariant functor}
    from $\calC\to \calD$ two maps:
    $F^{\text{obj}} : \operatorname{obj} \mathcal C
    \to\operatorname{obj} \mathcal D$
    and $F^{\text{Mor}} : 
    \operatorname{Mor}(A,B) \to 
    \boldsymbol{\operatorname{Mor}(F(B),F(A))}$
    satisfying similar axioms.
    The difference is that it reverses the direction of morphisms.)
\end{definition}
\begin{example}
    We have the following examples of functors.
    \begin{itemize}
	\item There is a functor from commutative rings
    to abelian groups that maps a ring $A$
    to a unit group $A^\times$.
    The morphism mapping is given by
    $(A\stackrel\phi\to B) \mapsto (A^\times\stackrel\phi\to B^\times)$
    by just restriction.
    \item There is a functor from commutative rings
	to commutative rings by $A\mapsto A[x]$.
    \end{itemize}
\end{example}
\section{Universal Properties}
\label{sec:universal}
\vocab[universal property]{Universal properties}
(also known as \vocab[universal mapping property]{universal mapping properties}
or \vocab[mapping property]{mapping properties})
determine an object up to isomorphism.
They are of a very specific format, which is best seen through examples.

\paragraph{Product}
Let $\calC$ be a category with objects $M$ and $N$.
Then, the product $M\times N$ is an object $P$, 
equipped with morphisms $\mu : P\to M$
and $\nu : P\to N$, such that the following universal property holds:
``for any object $P'$ and maps $\mu' : P'\to M$
and $\nu' : P'\to N$
\begin{center}
\begin{tikzcd}
    P' \arrow[rd, dashed] \arrow[rrd, "\mu'"]
    \arrow[rdd, "\nu'", swap]\\
    & P \arrow[r, "\mu", swap] \arrow[d, "\nu"] & M \\
    & N
\end{tikzcd},
\end{center}
there exists unique morphism $\phi : P'\to P$ 
(in dashed arrow) that makes the diagram above commute.''

(We say that a diagram of objects and morphisms \vocab{commutes}
if for any objects $P$ and $Q$ in the diagram,
every directed path from $P$ to $Q$ compose to the same 
map $P\to Q$.)

Let's suppose that $\calC$ is the category $\Sets$
(objects are sets, and morphisms are any function
between two sets). Given sets $M$ and $N$,
why does the definition above agree with the usual product
$$P = \{(m,n) : m\in M, n\in N\}?$$
First, the morphisms $\mu$ and $\nu$ are defined by
$\mu((m,n)) = m$ and $\nu((m,n)) = n$.
Now, let us verify that the above property holds.
Suppose that there exists $P'$, $\mu'$, and $\nu'$ as above.
Then, $\phi : P'\to P$ is uniquely determined to be 
$\phi(p) = (\mu'(p), \nu'(p))$.
Thus, $P$ satisfy the above universal property.

Now, we claim that \textbf{any two products satisfying
the universal property above are isomorphic}.
To prove this, suppose that we have two products $P_1$ and $P_2$:
\begin{center}
\begin{tikzcd}
    P_1 \arrow[r, "\mu_1", swap] \arrow[d, "\nu_1"] & M \\
     N
\end{tikzcd}\quad and\quad
\begin{tikzcd}
    P_2 \arrow[r, "\mu_2", swap] \arrow[d, "\nu_2"] & M \\
     N
\end{tikzcd}.
\end{center}
Then, by letting $P_2$ play the role of $P$
and $P_1$ play the role of $P'$ gives a morphism 
$\phi_1 : P_1\to P_2$ making the following diagram commute.
\begin{center}
\begin{tikzcd}
    P_2 \arrow[rd, bend left, "\phi_2"] 
    \arrow[rrd, bend left, "\mu_2"]
    \arrow[rdd, bend right, "\nu_2", swap]\\
    & P_1 \arrow[lu, bend left, "\phi_1", swap] 
    \arrow[r, "\mu_1", swap] \arrow[d, "\nu_1"] & M \\
    & N
\end{tikzcd}.
\end{center}
Similarly, there exists a morphism $\phi_2 : P_2\to P_1$
that makes the diagram above commute.

Now, we plug in $P=P'=P_1$. 
Then, there exists a unique morphism $P_1\to P_1$ 
that commutes with both $\mu_1$ and $\nu_1$.
However, we can name two such morphisms:
$\phi_2\circ\phi_1$ and the identity map.
Thus, $\phi_2\circ\phi_1 = \id_{P_1}$.
Similarly, $\phi_1\circ\phi_2 = \id_{P_2}$,
so $P_1$ and $P_2$ are isomorphic.

Notice how this argument does not use 
the nature of our category $\calC$ at all.
Thus, if one changes category $\calC$ to other categories
(e.g., $\Rings$ or $\Mod{R}$), then 
the exact same argument shows that any two products are isomorphic.
What's left, however, is to show that product
exists for each category that you want to define a product.

\paragraph{Polynomial Rings}
Let $R$ be a ring. Then, the polynomial ring $R[x]$
is the ring $S$ with element $x\in S$ and morphism 
$\iota : R\to S$ satisfying the following 
universal property: ``for ring $S'$, element $y\in S'$, 
and any morphism $\phi : R\to S'$,
there exists a unique morphism $f : S\to S'$
such that $f(x)=y$ and the diagram 
\begin{center}
\begin{tikzcd}
    & R \arrow[ld, "\iota", swap]  \arrow[rd, "\phi"] \\
    S \arrow[rr, "f"]& & S'  \\
    x \arrow[u, "\in", phantom, sloped]
    \arrow[rr, maps to] & & y \arrow[u, "\in", phantom, sloped] 
\end{tikzcd}
\end{center}
commutes".
\begin{exercise}
Using the same argument as products, 
verify that the universal property given above 
characterizes $S$ up to isomorphism.
\end{exercise}
We now check that a ring $S$ satisfying the above universal 
property exists. 
To do this, we give the usual construction: define 
$$S = \{a_0 + a_1x + \dots + a_nx^n : 
a_0,a_1,\dots,a_n\in R\},$$
and the addition / multiplication operations 
are the standard polynomial addition / multiplication.
Then, for any $S'$ and $\phi$ satisfying the above,
one can see that the map $f : S\to S'$ is forced to be
$$f(a_0 + a_1x + \dots + a_nx^n)
= \phi(a_0) + \phi(a_1)y + \dots + \phi(a_n)y^n.$$
We leave the reader to check that this defines 
a valid morphism.
To do this, one needs to check that it is compatible 
under polynomial addition (easy) and multiplication
(harder, but still straightforward).
\paragraph{Universal Property in General}
In general, the recipe for defining an object through 
universal properties is as follows.
\begin{enumerate}
    \item Define the universal property.
    \item Prove that this property characterizes an object 
    up to an isomorphism (a routine process).
    \item Give a construction of an object
    (possibly difficult).
\end{enumerate}
In these notes, we will soon see examples (e.g., tensor product, \Cref{sec:tensor}) where Step 3.
is nontrivial, and therefore, the construction of the object
itself makes it very hard to prove properties.
When such situation arises, we often forget about Step 3.
and use universal properties to deduce properties about that object.
\begin{remark}[Towards Yoneda's lemma]
In fact, one can formalize a theorem which says that 
``any universal property determines an object up to isomorphism''.
The result is \vocab{Yoneda's lemma},
which require a little bit more category theory to state
than we want to do here (see \cite[\S 4.2]{category} for more details).

As a starting point towards Yoneda's lemma,
let us rephrase the universal properties of product 
and polynomial rings a little bit.
\begin{itemize}
\item $(P,\mu, \nu)$ is a product if for each $P'$,
there exists a bijective corresondence between
$$\{\text{morphisms } \phi : P'\to P\} 
\ \longleftrightarrow\ 
\{(\text{morphisms }\mu' : P'\to M \text{ and } \nu' : P'\to N)\}.$$
\item $(S,x)$ is a polynomial ring if for each $S'$,
there exists a bijective correspondence between 
$$\{\text{morphisms } S\to S'\} 
\ \longleftrightarrow\ 
\{\text{morphisms } R\to S', \text{ image of }x\}.$$
\end{itemize}
In the first case, the universal property is classifying 
morphisms from any object to $P$,
while in the second case, the universal property is classifying
morphisms from $S$ to any object.

These represent two types of universal properties of $X$:
properties classifying the set of morphisms \textbf{from} $X$
and properties classifying the set of morphisms \textbf{to} $X$.
Yoneda's lemma states that, provided with some other mild conditions,
properties of these two types classify objects up to isomorphism.
\end{remark}
\section{Topological Spaces}
\label{sec:topological_space}
We give a quick review of topological spaces
and topological notions used in commutative algebra.
\begin{definition}
A \vocab{topological space} $(X,\calT)$ consists of 
\begin{itemize}
    \item a set $X$ of points.
    \item a set $\calT\subseteq 2^X$ of \vocab{open sets}
\end{itemize}
that satisfy the following conditions:
\begin{itemize}
    \item $\emptyset$ and $X$ are both in $\calT$.
    \item for an arbitrary (possibly infinite, even uncountable) collection 
    $(U_i)_{i\in I}$ where $U_i\in\calT$,
     we have $\bigcup_{i\in I} U_i\in \calT$.
    \item for any $U_1,U_2\in\calT$, we have 
    $U_1\cap U_2\in\calT$.
\end{itemize}
A set is \vocab[closed sets]{closed} if and only if its complement is open.
\end{definition}
Originally, these terms are constructed to model geometric 
spaces such as $\bbR^n$ or a torus (see \cite{topology} 
for more details). 
However, in algebra, we often borrow words from topology to describe structure that may or may not have geometric meaning.
The words necessary for commutative algebra are as follows.
\begin{definition}
    Let $(X,\calT)$ be a topological space,
    and let $S\subseteq X$.
    Then, the \vocab{induced topology}
    on $S$ is defined by the list of open sets
    $$\calT' = \{V\subseteq S : V=U\cap S
    \text{ for some open set }U\in\calT\}.$$
\end{definition}
\begin{definition}[Bases]
    Let $(X,\calT)$ be a topological space.
    A collection of open sets $\calB\subseteq \calT$
    forms a \vocab{bases} if and only if any open set 
    in $\calT$ can be written as a union of sets in $\calB$.
\end{definition}

\begin{definition}[Connectedness and Quasicompactness]
\label{def:compact}
A topological space $(X,\calT)$ is said to be 
\begin{itemize}
    \item \vocab{connected} if $\emptyset$ and $X$
    are the only subsets that are both closed and open
    (equivalently, $X$ cannot be partition into 
    disjoint union of open sets $U_1\sqcup U_2$).
    \item \vocab{quasicompact} if any open covering has a finite subcovering. In other words, for any collection $(U_i)_{i\in I}$ such that $\bigcup_{i\in I}U_i=X$,
    there exists a finite subset $J\subseteq I$ such that 
    $\bigcup_{j\in J} U_j = X$.
\end{itemize}
\end{definition}
\begin{remark}
    In some other contexts, the definition of quasicompact 
    given above is the definition of compact,
    but in some places, compact means quasicompact and Hausdorff. Thus, we use the word quasicompact to make clear that we do not need Hausdorff condition.
\end{remark}
\begin{exercise}
    Prove that any closed subset (with its induced topology) of a quasicompact space is quasicompact.
\end{exercise}
